\scanpage{455}\mspage{367}
% archive sheet number: 454
$q\in\Sigma_0$ tel que $qp=qp'$ (i.e telle que la sous-catégorie \inserted{$J_0$} de $J$ dont l'ens des objets est $\{e\}$, et l'ens des flèches $\{\mathrm{id}_e\}\cup\Sigma_{\infty}$, soit filtrante).

Si $\Sigma_0$ est choisie ainsi, soit $F$ l'ens \deleted{des} \inserted{réunion de $\mathrm{Fl}\,\deleted{\unsure{J}}K$ et de l'ens des} flèches \deleted{de la \unsure{fo-}} composées
\diagram{\begin{tikzcd}[ampersand replacement=\&]
\unsure{(\alpha)|(a)} \& k\arrow[r,"\unsure{u}"] \& e\arrow[r,"p"] \& \deleted{\unsure{k}}e\arrow[r,"v"] \& k',
\end{tikzcd}}
avec $u,v\in\mathrm{Fl}\,J$, $p\in\Sigma_{\infty}$. Alors $F$ est l'ens. des flèches d'une \inserted{sous-}catégorie $\deleted{\unsure{J'}}K'$ de $J$, ayant \unsure{mes|m\^es|mêmes} objets que $J$, \deleted{et $K'$ est \unsure{fi-}} et admettant $e$ pour plus grand \unsure{élémt}. De plus, $K'$ contient $K$ et est filtrante.

\uline{\unsure{Dém}} Que $\Sigma$ soit stable par composition est clair. Pour construire $\Sigma_{\infty}$ ayant les propriétés énoncées, on note que, $\deleted{K}\unsure{J|K}$ \deleted{\unsure{étant}} \unsure{gde|grde} devant $\pi$, \deleted{et \unsure{f-} il existe} $\Sigma\deleted{\unsure{_0}}$ est non vide.$^*$
% slanted
\margin{cf argument plus bas pour \unsure{constr.|const.} de $p_i$} On construit \inserted{par \unsure{ordre|ind.}} une suite d'élémts
\[
p_0,p_1,\ldots\ p_n,\ldots
\]
de $\Sigma$, en prenant pour $p_0$ un élément quelconque de $\Sigma$, et pour $\unsure{i|j}\geq1$, prenant $p_i\in\Sigma$ tel que \inserted{pour} \unsure{tte} \deleted{\unsure{double}} double-flèche \deleted{\unsure{$(\alpha,\beta)$}} $(p,p')$ du \inserted{\unsure{ss-|s-}}monoïde \inserted{$M_{i-1}$} de $\mathrm{End}_J(e)$ engendré par $(p_j)_{0\leq j\leq i-1}$, \deleted{on} ait $p_i p=p_i p'$. Un tel $p_i$ existe, \deleted{puisque \unsure{par} \inserted{on \deleted{\unsure{remarque|puisque}}} \unsure{l'hyp.} \unsure{de} $J$ est \unsure{grand|grd} devant $\pi$ \unsure{donc} devant $\aleph_0$, et que $M_{i-1}$ est de cardinal $\leq\aleph_0$, il} en prenant l'ens des doubles-flèches de la forme précédente $(p,p')$, avec $p,p'\in M_{i-1}$, et de celles de la forme $(u,v):k\rightrightarrows e$, avec $u,v\in\mathrm{Fl}\,J$, l'ens des
